\documentclass{article}

\usepackage{graphicx}
\usepackage{amsmath}
\usepackage[top=0.8in, bottom=1in, left=0.8in, right=0.8in]{geometry}
\usepackage{titlesec}
\usepackage{float}
\usepackage{booktabs}
\usepackage{subcaption}
\usepackage{listings}
\usepackage{xcolor}

\lstdefinestyle{sqlstyle}{
    language=SQL,
    basicstyle=\ttfamily\small,
    keywordstyle=\color{blue},
    stringstyle=\color{red},
    commentstyle=\color{gray},
    numbers=left,
    numberstyle=\tiny\color{gray},
    stepnumber=1,
    numbersep=8pt,
    frame=single,
    breaklines=true,
    showstringspaces=false
}

\titlespacing*{\section}{0pt}{1em}{0.5em}
\titlespacing*{\subsection}{0pt}{0.5em}{0.3em}

\setlength{\parindent}{0pt}
\setlength{\parskip}{0.4em}

\usepackage{float}

\usepackage{fancyhdr}
\pagestyle{fancy}

\fancyhf{}
\fancyhead[C]{MIT Sloan School of Management --- 15.458 Financial Data Science and Computing}
\fancyfoot[C]{\thepage}

\renewcommand{\headrulewidth}{0.4 pt}

\title{\textbf{Decoding Insider Trading Signals}\\
\bigskip 

\large \textbf{Routine vs.\ Opportunistic Trades, Machine Learning, and Return Predictability}}
\author{Shabbir Houssenaly \and Camilo Sanchez Quinto \and Brandon Li \and Niccol\`o Nordera}

\begin{document}
\maketitle

%------------------------------------------------------------------
\begin{abstract}
\noindent We study whether insider trades predict future stock returns using 1,716,763 SEC Form 4 transactions across 8,601 U.S. firms from 2010 to 2025. We classify each trade as routine or opportunistic following Cohen, Malloy, and Pomorski (2012), applying the criterion year-by-year rather than permanently. Insider signals are aggregated using SEC filing dates rather than transaction dates, ensuring that only publicly available information enters the analysis. Panel regressions show that opportunistic insider buying predicts next-month returns of 0.35\% ($p < 0.001$), with all predictive content concentrated in high-volatility environments. A gradient boosting model (XGBoost) trained on an expanding window generates 735,365 out-of-sample predictions; a long-short portfolio earns 34.99\% annualized (Sharpe 1.60), growing \$1 to \$12.86, with a Fama-French five-factor alpha of 2.78\% per month ($t = 3.91$). A Proximal Policy Optimization (PPO) reinforcement learning agent that directly optimizes the Sharpe ratio with explicit transaction cost penalties earns 74.18\% annualized over 2020--2024 (Sharpe 0.91). Feature importance reveals that market and fundamental variables drive ML predictability, while the RL agent's higher raw returns reflect its ability to take concentrated positions in high-conviction months.
\end{abstract}

%------------------------------------------------------------------
\section{Introduction}
%------------------------------------------------------------------

Corporate insiders possess private information not yet reflected in prices. When they trade in their own firm's securities, they implicitly signal their private assessment of firm value. However, not all insider trades are equally informative: many reflect liquidity needs, diversification, or pre-scheduled Rule 10b5-1 plans. Cohen, Malloy, and Pomorski (2012) distinguish \textit{routine} trades --- those occurring in the same calendar month for at least three consecutive years --- from \textit{opportunistic} trades that deviate from any established pattern, showing that only opportunistic trades predict future returns.

We extend this framework in four directions. First, we apply the routine criterion year-by-year rather than permanently, reducing routine trades from $\sim$55\% to 27\%. Second, we aggregate insider signals using SEC filing dates rather than transaction dates, ensuring strict informational correctness: the signal in month $t$ reflects only trades whose Form 4 was publicly filed by month-end. Third, we train XGBoost on an expanding window and construct long-short portfolios with large out-of-sample alpha confirmed by Fama-French five-factor regressions. Fourth, and distinctively, we apply a PPO reinforcement learning agent that directly optimizes the portfolio Sharpe ratio with explicit transaction cost penalties.

%------------------------------------------------------------------
\section{Data}
%------------------------------------------------------------------

We combine SEC Form 4 insider transactions (WRDS), monthly CRSP stock returns, and quarterly Compustat fundamentals into a firm-month panel. Starting from 1,742,812 transactions (2010--2025), we retain Form 4 open-market purchases and sales with non-missing prices, yielding \textbf{1,716,763 transactions}. We link firms to CRSP via CIK$\to$GVKEY$\to$PERMNO, matching 8,601 firms (75.8\%), and resolve 211,920 duplicate firm-month observations from many-to-many CIK-PERMNO links. The final panel contains \textbf{1,110,501 firm-month observations} over January 2010 -- December 2025.

A methodological refinement relative to prior work is that we aggregate insider signals using the SEC \textit{filing date} (\texttt{fdate}) rather than the \textit{transaction date}. Form 4 must be filed within two business days of the transaction, but the trade becomes public information only upon filing. Using filing dates ensures that the signal available at month $t$ reflects only trades whose Form 4 was publicly disclosed by month-end, avoiding any look-ahead bias from transaction timing.

At the transaction level, we retain open-market purchases and sales and construct signed trade measures using purchases as positive and sales as negative. We first aggregate transactions to the insider--firm--filing-date level and then to the firm--filing-month level. This aggregation captures both the direction and intensity of insider activity while preserving the public-information timing of the filing. The resulting insider variables include net signed trade value, gross purchase and sale values, gross purchase and sale shares, number of transactions, number of unique insiders, officer and director participation, ownership-scaled trade size, and a coordination dummy equal to one when multiple insiders trade in the same firm-month. We then merge these insider signals to CRSP returns and Compustat fundamentals through the CIK$\to$GVKEY$\to$PERMNO link.

Control variables include momentum, trailing 12-month volatility, book-to-market, gross profitability, and firm size. Momentum is defined as the compounded stock return from months $t$-12 to $t$-2, skipping the most recent month. Volatility is the trailing 12-month standard deviation of monthly returns. Book-to-market is constructed from Compustat book equity and CRSP market equity using standard timing conventions, while firm size is measured as the logarithm of CRSP market equity. Continuous variables are winsorized at the 1st/99th percentiles and accounting variables are forward-filled to monthly frequency.

\begin{table}[h!]
\centering
\caption{Variable Construction}
\label{tab:variables}
\small
\begin{tabular}{p{0.24\linewidth}p{0.22\linewidth}p{0.45\linewidth}}
\toprule
Variable & Source & Construction \\
\midrule
Return ($r_{t+1}$) & CRSP & One-month-ahead stock return used as the dependent variable. \\
Momentum & CRSP & Compounded return from months $t$-12 to $t$-2. \\
Volatility & CRSP & Trailing 12-month standard deviation of monthly returns. \\
Firm Size & CRSP & Log of market equity, where market equity equals absolute price times shares outstanding. \\
Book-to-Market & Compustat/CRSP & Book equity divided by market equity, aligned using standard accounting-data timing conventions. \\
Gross Profitability & Compustat & Gross profitability measure forward-filled to monthly frequency. \\
Net Insider Trading & Form 4 & Signed net trade value, with purchases positive and sales negative. \\
Gross Buy/Sell Value & Form 4 & Total dollar value of purchases and sales within a firm-month. \\
Trade Direction & Form 4 & Indicator equal to $+1$ for net buying, $-1$ for net selling, and $0$ for mixed or zero net activity. \\
Insider Role & Form 4 & Officer and director indicators, including officer-title information when available. \\
Ownership-Scaled Trade Size & Form 4 & Transaction shares scaled by post-transaction ownership. \\
Coordination Dummy & Form 4 & Indicator equal to one when at least two distinct insiders trade in the same firm-month. \\
Routine/Opportunistic Signal & Form 4 & Insider trades classified using the Cohen, Malloy, and Pomorski timing rule and aggregated to the firm-month level. \\
\bottomrule
\end{tabular}
\begin{minipage}{\linewidth}
\vspace{2pt}
\small\textit{Notes:} Insider variables are aggregated using SEC filing dates rather than transaction dates to ensure that all signals are publicly observable before portfolio formation or return prediction.
\end{minipage}
\end{table}

\begin{table}[H]
\centering
\caption{Summary Statistics}
\label{tab:summary}
\small
\begin{tabular}{lrrrrr}
\toprule
Variable & N & Mean & Std Dev & P25 & P75 \\
\midrule
Return ($r_{t+1}$)    & 1,076,147 & 0.004 & 0.095 & $-$0.056 & 0.058 \\
Volatility            & 1,105,183 & 0.120 & 0.083 & 0.059    & 0.153 \\
Momentum              & 1,112,826 & 0.066 & 0.428 & $-$0.177 & 0.240 \\
Book-to-Market        &   878,083 & 0.561 & 0.712 & 0.206    & 0.866 \\
Gross Profitability   &   808,390 & 0.059 & 0.072 & 0.014    & 0.099 \\
Signal\_Opportunistic & 1,110,501 & 0.054 & 0.234 & 0        & 0     \\
Signal\_Routine       & 1,110,501 & 0.056 & 0.247 & 0        & 0     \\
\bottomrule
\end{tabular}
\begin{minipage}{\linewidth}
\vspace{2pt}
\small\textit{Notes:} Firm-month panel, January 2010 -- December 2025. All continuous variables winsorized at the 1st/99th percentiles. Insider signals aggregated by SEC filing date.
\end{minipage}
\end{table}

%------------------------------------------------------------------
\section{Classification}
%------------------------------------------------------------------

Following Cohen et al.\ (2012), a trade in year $t$ is classified as \textit{routine} if the insider traded in the same calendar month in each of the two preceding years. All other trades are \textit{opportunistic}. Our year-by-year implementation ensures that once an insider breaks a consecutive streak, subsequent trades immediately revert to opportunistic --- a stricter application that reduces the routine share to 27\% (vs.\ $\sim$55\% in a permanent-label implementation). Directional signals $\text{Signal}_{i,t} \in \{-1, 0, +1\}$ capture net buying direction at the firm-month level, winsorized at the 1st/99th percentiles.

\begin{table}[h!]
\centering
\caption{Routine vs.\ Opportunistic Classification}
\label{tab:classification}
\small
\begin{tabular}{lrr}
\toprule
Category & Transactions & Share \\
\midrule
Routine       & 464,167   & 27.0\% \\
Opportunistic & 1,252,596 & 73.0\% \\
Total         & 1,716,763 & 100\%  \\
\bottomrule
\end{tabular}
\begin{minipage}{\linewidth}
\vspace{2pt}
\small\textit{Notes:} Year-by-year application reduces routine share to 27\% vs.\ $\sim$55\% in a permanent-label implementation.
\end{minipage}
\end{table}

%------------------------------------------------------------------
\section{Panel Regression Results}
%------------------------------------------------------------------

We estimate three panel regressions with \texttt{feols} (R \texttt{fixest}), using firm and time fixed effects and firm-clustered standard errors on 600,699 firm-month observations. The dependent variable is $R_{i,t+1}$.

\vspace{4pt}
\noindent\textbf{Model 1} (baseline): $R_{i,t+1} = \alpha_i + \gamma_t + \beta_1 \text{SigOpp}_{i,t} + \beta_2 \text{SigRou}_{i,t} + \delta' X_{i,t} + \varepsilon_{i,t}$

\noindent\textbf{Model 2} (dummy): $R_{i,t+1} = \alpha_i + \gamma_t + \beta_1 \mathbf{1}[\text{Opp}_{i,t}>0] + \delta' X_{i,t} + \varepsilon_{i,t}$

\noindent\textbf{Model 3} (interaction): adds $\beta_3(\text{SigOpp}_{i,t} \times \text{Vol}_{i,t})$ to Model 1.

\vspace{4pt}
\noindent A Ljung-Box test (lag 12) on firm-level residuals finds that 9.6\% of 6,269 firms reject white noise at 5\%, marginally above the 5\% false-positive rate. Firm-clustered standard errors are adequate.

\begin{table}[h!]
\centering
\caption{Panel Regression Results}
\label{tab:regressions}
\small
\begin{tabular}{lccc}
\toprule
 & Model 1 & Model 2 & Model 3 \\
\midrule
Signal\_Opportunistic           & $0.0035^{***}$ &                 & $0.0009$       \\
                                & $(0.0004)$     &                 & $(0.0007)$     \\
Signal\_Routine                 & $0.0044^{***}$ &                 & $0.0043^{***}$ \\
                                & $(0.0005)$     &                 & $(0.0005)$     \\
Any\_Opportunistic              &                & $-0.0021^{***}$ &                \\
                                &                & $(0.0004)$      &                \\
SigOpp $\times$ Volatility      &                &                 & $0.0203^{***}$ \\
                                &                &                 & $(0.0059)$     \\
Momentum                        & $-0.0042^{***}$ & $-0.0045^{***}$ & $-0.0041^{***}$ \\
Volatility                      & $0.0114^{*}$   & $0.0124^{**}$  & $0.0123^{**}$  \\
Book-to-Market                  & $0.0128^{***}$ & $0.0129^{***}$ & $0.0128^{***}$ \\
Gross Profitability             & $0.1335^{***}$ & $0.1329^{***}$ & $0.1337^{***}$ \\
\midrule
Firm FE / Time FE  & Yes/Yes   & Yes/Yes   & Yes/Yes   \\
Observations       & 600,699   & 600,699   & 600,699   \\
Within R$^2$       & 0.0041    & 0.0039    & 0.0041    \\
\bottomrule
\end{tabular}
\begin{minipage}{\linewidth}
\vspace{2pt}
\small\textit{Notes:} SE clustered by firm. Controls shown without SE for space. Signals aggregated by SEC filing date. $^{***}p<0.001$, $^{**}p<0.01$, $^{*}p<0.05$.
\end{minipage}
\end{table}

\paragraph{Key findings.} Opportunistic buying predicts 0.43\% higher next-month returns ($p<0.001$), or $\sim$5.2\% annualized. 
The positive coefficient on routine trades is unexpected under the strictest interpretation
of the routine-trading hypothesis. One possibility is that our year-by-year classification
identifies a subset of persistent insiders whose trades still contain information. Another
possibility is that the directional routine signal captures residual firm characteristics or
unmodeled trade-size/role effects. We therefore interpret the opportunistic/routine contrast
with caution.
A binary opportunistic indicator is negative ($-$0.0013), confirming that \textit{direction}, not occurrence,  drives predictability. In Model 3, the direct opportunistic effect vanishes while the interaction with volatility is 0.0296 ($p<0.001$): the predictive content of opportunistic trades is substantially stronger in high-volatility environments, consistent with insiders having greater informational advantage when public information is sparse. 

%------------------------------------------------------------------
\section{Machine Learning, Portfolio Performance, and Alpha}
%------------------------------------------------------------------

\subsection{XGBoost with Expanding Window}

We train a gradient boosting classifier (XGBoost, \texttt{binary:logistic}) to predict whether firm $i$ earns a positive return in month $t+1$:
\[
    y_{i,t+1} = \mathbf{1}\{R_{i,t+1} > 0\}.
\]
The model outputs a predicted probability $\hat{p}_{i,t}$ for each firm-month. We estimate the model using an expanding-window design: for each test year $t \in \{2015,\ldots,2024\}$, the model is trained only on data from 2010 through $t-1$ and then used to generate out-of-sample predictions for year $t$. Features include 19 variables spanning insider signals, trade counts and volumes, market variables, and accounting fundamentals. Early stopping with 30 rounds and a 10\% validation holdout limits overfitting, yielding \textbf{735,365 out-of-sample firm-month predictions}.

We convert predicted probabilities into trading signals by sorting firms by $\hat{p}_{i,t}$ each month. The main XGBoost portfolio goes long firms in the top decile and short firms in the bottom decile. Returns are equal-weighted within each leg and held for one month.

Figure \ref{fig:importance} shows that stock price (33.1\%) and volatility (25.0\%) dominate predictive gain, with all insider features contributing below 2\% total. XGBoost therefore outperforms the linear insider-signal benchmark mainly by capturing nonlinear interactions among market and fundamental predictors, consistent with Gu, Kelly, and Xiu (2020).

\begin{figure}[h!]
\centering
\includegraphics[width=0.76\textwidth]{fig_importance.pdf}
\caption{XGBoost feature importance (Gain), model trained on 2010--2023. Purple = market variables, green = fundamentals, coral = insider signals.}
\label{fig:importance}
\end{figure}

\subsection{Reinforcement Learning Portfolio Optimization}

To directly optimize the portfolio Sharpe ratio while accounting for transaction costs, we apply a Proximal Policy Optimization (PPO) reinforcement learning agent (Schulman et al., 2017). Unlike XGBoost, which predicts returns for each stock independently and applies a fixed trading rule, the RL agent learns an adaptive allocation policy across the 10 XGBoost probability deciles.

\paragraph{Environment design.} At each month $t$, the agent observes a state vector of 21 dimensions: the mean XGBoost probability and mean opportunistic signal for each of the 10 deciles (20 features), plus the previous month's portfolio return. The action space is continuous: the agent outputs a weight $w_k \in [-1, +1]$ for each decile $k$.

\paragraph{Reward function.} The monthly reward is the rolling 12-month Sharpe ratio of portfolio returns, net of transaction costs:
\[
r_t = \frac{\bar{R}_{t-11:t}}{\sigma_{t-11:t} + \epsilon} - \lambda \sum_{k=1}^{10} |w_{k,t} - w_{k,t-1}|
\]
where $\lambda = 0.001$ (10 basis points per unit of turnover) penalizes excessive trading.

\paragraph{Architecture and training.} The policy network is a three-layer MLP (128-128-64 neurons, ReLU activations, Tanh output). We use an expanding window identical to XGBoost: the agent trains from scratch on data from 2015 through $t-1$ before predicting year $t$, for $t \in \{2020, \ldots, 2024\}$. The first test year is 2020, as the agent requires at least 5 years of decile-level training data.

\subsection{Portfolio Performance Comparison}

Table \ref{tab:portfolio} reports performance statistics and Figure \ref{fig:wealth} plots cumulative wealth.

\begin{table}[h!]
\centering
\caption{Portfolio Performance}
\label{tab:portfolio}
\small
\begin{tabular}{lcccccc}
\toprule
Strategy & Period & Ann.\ Return & Volatility & Sharpe & Hit Rate & \$1 $\to$ \\
\midrule
Regression L/S       & 2015--2024 & $-2.55\%$  & $6.64\%$  & $-0.38$ & $44.3\%$ & \$0.78  \\
XGBoost Top/Bot 20\% & 2015--2024 & $+25.27\%$ & $18.23\%$ & $+1.39$ & $71.4\%$ & \$6.69  \\
XGBoost Decile 10-1  & 2015--2024 & $+34.99\%$ & $21.88\%$ & $+1.60$ & $73.9\%$ & \$12.86 \\
RL PPO (continuous)  & 2020--2024 & $+74.18\%$ & $81.53\%$ & $+0.91$ & $62.7\%$ & \$4.65  \\
\bottomrule
\end{tabular}
\begin{minipage}{\linewidth}
\vspace{2pt}
\small\textit{Notes:} Equal-weighted within legs for regression and XGBoost strategies; continuous weights for RL. Monthly rebalancing, gross of transaction costs for regression and XGBoost; net of 10 bps turnover penalty for RL. RL test period starts in 2020 due to training data requirements. All signals aggregated by SEC filing date.
\end{minipage}
\end{table}

\begin{figure}[h!]
\centering
\includegraphics[width=0.76\textwidth]{fig_wealth.pdf}
\caption{Cumulative wealth of \$1 invested at strategy inception, monthly rebalancing. XGBoost Decile (solid purple, 2015--2024), XGBoost 20\% (dashed green, 2015--2024), Regression L/S (dashed gray, 2015--2024). RL PPO results reported separately in Table \ref{tab:portfolio} due to different test period.}
\label{fig:wealth}
\end{figure}

The regression-based strategy underperforms ($-2.55\%$, Sharpe $-0.38$) when using filing dates, illustrating both the gap between in-sample statistical significance and out-of-sample economic value, and the slight attenuation from using strictly public information. XGBoost strategies generate substantially larger returns, with the decile strategy achieving Sharpe 1.60 and growing \$1 to \$12.86.

The RL agent generates the highest raw return (+74.18\% annualized) but with substantially higher volatility (81.5\%), resulting in a Sharpe ratio of 0.91. Year-by-year, the agent performs strongly in 2021 (Sharpe 2.37) and 2024 (Sharpe 1.98) but is more volatile in other years. On a risk-adjusted basis, XGBoost remains the dominant strategy.

\subsection{Fama-French Five-Factor Alpha}

We regress each strategy's monthly returns on the Fama-French five factors using Newey-West standard errors (lag = 6). Table \ref{tab:ff5} reports results.

\begin{table}[h!]
\centering
\caption{Fama-French 5-Factor Alpha (Filing-Date Signals)}
\label{tab:ff5}
\small
\begin{tabular}{lccccc}
\toprule
Strategy & $\alpha$/month & $\alpha$/year & $t$-stat & $R^2$ & $N$ \\
\midrule
Regression L/S       & $-0.0023$       & $-2.77\%$  & $-1.50$ & 0.020 & 106 \\
XGBoost Top/Bot 20\% & $+0.0203^{***}$ & $+24.41\%$ & $+3.68$ & 0.050 & 119 \\
XGBoost Decile 10-1  & $+0.0278^{***}$ & $+33.32\%$ & $+3.91$ & 0.048 & 119 \\
\bottomrule
\end{tabular}
\begin{minipage}{\linewidth}
\vspace{2pt}
\small\textit{Notes:} Newey-West SE, lag=6. FF5 factors from Kenneth French Data Library. RL not included due to short sample (59 months). $^{***}p<0.001$.
\end{minipage}
\end{table}

The regression strategy generates no significant alpha ($-2.77\%$, $p=0.14$). Both XGBoost strategies generate large, statistically significant alphas: $+24.41\%$ per year ($t=3.68$) and $+33.32\%$ per year ($t=3.91$). The low $R^2$ (below 6\%) indicates that FF5 factors explain little of the portfolio return variation. Notably, the XGBoost alpha is essentially unchanged relative to transaction-date results ($+32.61\%$ vs $+33.32\%$), confirming robustness to the filing-date refinement.

%------------------------------------------------------------------
\section{Conclusion}
%------------------------------------------------------------------

This paper studies the predictability of insider trades using panel regressions, gradient boosting, and reinforcement learning. A key methodological contribution is the use of SEC filing dates rather than transaction dates for signal aggregation, ensuring strict informational correctness. Our main findings are: (i) opportunistic insider trades predict next-month returns of 0.35\% using filing-date signals, concentrated entirely in high-volatility environments; (ii) XGBoost generates large out-of-sample alpha (33.32\% per year, FF5-adjusted) robust to the filing-date refinement, by capturing non-linear combinations of market and fundamental variables; (iii) a PPO reinforcement learning agent generates 74.18\% annualized returns (Sharpe 0.91) by directly optimizing the Sharpe ratio with transaction cost penalties, though at substantially higher volatility than XGBoost.

Several limitations should be noted. All portfolio returns are reported gross of transaction costs, which could materially reduce net performance. The RL test period covers only 2020--2024 and includes unusual market conditions. Future work could exploit post-2023 Form 4 plan indicators to separately identify Rule 10b5-1 plan trades, apply Fama-French decomposition to RL returns over a longer sample, and incorporate value-weighted portfolio construction.


%------------------------------------------------------------------
\appendix
\section*{Appendix: Reinforcement Learning Design Details}
%------------------------------------------------------------------

\paragraph{State space.} The 21-dimensional state vector at time $t$ consists of: (1) mean XGBoost predicted probability for each of the 10 deciles (10 features); (2) mean \texttt{signal\_opportunistic} for each decile (10 features); (3) the previous month's portfolio return (1 feature).

\paragraph{Action normalization.} Raw actions from the policy network $a \in [-1,+1]^{10}$ are normalized so the maximum absolute weight equals 1.

\paragraph{Hyperparameters.} Learning rate $3 \times 10^{-4}$; $n$-steps 512; batch size 64; 10 epochs per update; GAE lambda 0.95; clip range 0.2; entropy coefficient 0.01. Policy network: MLP with layers [128, 128, 64], ReLU activations, Tanh output layer.

\paragraph{Expanding window.} For test year $t$, the agent trains on monthly decile data from 2015 through $t-1$. The first test year is 2020 (60 training months, 30,000 training steps). Training steps scale with the training window size ($500 \times$ number of training months).

\paragraph{Transaction cost model.} The turnover penalty $\lambda \sum_k |w_{k,t} - w_{k,t-1}|$ with $\lambda = 0.001$ approximates 10 basis points per unit of portfolio weight changed.

%------------------------------------------------------------------
\section*{References}
%------------------------------------------------------------------
\small
\noindent Chen \& Guestrin (2016). XGBoost. \textit{KDD}, 785--794. \quad
Cohen, Malloy \& Pomorski (2012). Decoding inside information. \textit{JF}, 67(3), 1009--1043.

\noindent Fama \& French (1993). Common risk factors. \textit{JFE}, 33(1), 3--56. \quad
Fama \& French (2015). A five-factor asset pricing model. \textit{JFE}, 116(1), 1--22.

\noindent Freyberger, Neuhierl \& Weber (2020). Dissecting characteristics. \textit{RFS}, 33(5), 2326--2377. \quad
Gu, Kelly \& Xiu (2020). Empirical asset pricing via ML. \textit{RFS}, 33(5), 2223--2273.

\noindent Jaffe (1974). Special information and insider trading. \textit{JB}, 47(3), 410--428. \quad
Schulman et al.\ (2017). Proximal policy optimization algorithms. \textit{arXiv:1707.06347}.

\noindent Seyhun (1986). Insiders' profits. \textit{JFE}, 16(2), 189--212.

\end{document}


